% (C) 2026 Ralf Stephan, in collaboration with Claude Code.
% Released under CC0 1.0 Universal (public-domain dedication).
% See https://creativecommons.org/publicdomain/zero/1.0/

\documentclass[11pt]{article}
\usepackage[T1]{fontenc}
\usepackage[utf8]{inputenc}
\usepackage{lmodern}
\usepackage[margin=2.5cm]{geometry}
\usepackage{amsmath,amssymb}
\usepackage{booktabs,tabularx,xltabular}
\usepackage{braket}
\usepackage[hidelinks]{hyperref}

% Lean names and paths: verbatim, breakable at . _ /
\DeclareUrlCommand\lean{\urlstyle{tt}}
\def\UrlBreaks{\do\.\do\_\do\/}
\newcolumntype{L}{>{\raggedright\arraybackslash}X}
\newcolumntype{B}{>{\raggedright\arraybackslash}p{0.42\textwidth}}
\emergencystretch=6em

\title{Lean formalization of an Assumptions of Physics technical brief:\\ log}
\author{}
\date{}

\begin{document}
\maketitle

\noindent
This log records which statements of the technical brief in this directory were checked in Lean,
where the formal version differs from the brief, and any problems found in the brief.
The brief is available at \url{https://assumptionsofphysics.org/resources/briefs}.

\section*{Setup}

\begin{itemize}
\item Lean \texttt{v4.36.0-rc2}; Mathlib at commit \texttt{1b1ccf4b}; physlib (\texttt{master})
  and TauCeti (commit \texttt{e7e7f52}, on the same Mathlib) as further dependencies.
\item Build: \texttt{lake build} in the project root.
\item Every file is checked with \texttt{\#print axioms} on its main results. A result is
  recorded here only if it depends on no axioms beyond \lean{propext},
  \lean{Classical.choice} and \lean{Quot.sound}, and so has no \lean{sorry}.
\item Imports are kept to the Mathlib modules actually needed.
\item Material that belongs in Mathlib is kept in \lean{ForMathlib/}, in Mathlib's directory
  layout (library \lean{ForMathlib}), and imported by the brief files.
\end{itemize}


%======================================================================

\section{Classical uncertainty principle}

\paragraph{Files.}\leavevmode
\begin{itemize}
\item File: \lean{briefs/ClassicalUncertaintyPrinciple/ClassicalUncertaintyPrinciple.lean}
\item Library: \lean{ClassicalUncertaintyPrinciple}
\item Imports: six new \lean{ForMathlib} files, the first two through
  \lean{Subadditivity.lean}:
  \begin{itemize}
  \item \lean{ForMathlib/InformationTheory/GibbsInequality.lean}:\newline Gibbs' inequality
  $-\int\rho\log\rho \le -\int\rho\log G$ for a probability density $\rho$ and a positive $G$
  with $\int G \le 1$ (\lean{MeasureTheory.integral_negMulLog_le}), with equality only for
  $\rho = G$ a.e.\ (\lean{MeasureTheory.ae_eq_of_integral_negMulLog_eq}); the same for $G \ge 0$
  positive wherever $\rho$ is (\lean{MeasureTheory.integral_negMulLog_le_of_imp_pos})
  \item \lean{ForMathlib/Probability/Distributions/Gaussian/Real.lean}: the logarithm of the
  Gaussian density (\lean{log_gaussianPDFReal}), and its mean and variance as Lebesgue integrals
  (\lean{integral_mul_gaussianPDFReal}, \lean{integral_sub_sq_mul_gaussianPDFReal})
  \item \lean{ForMathlib/InformationTheory/EntropyVariation.lean}:\newline the first variation
  $\frac{d}{dt}\int -\rho_t\log\rho_t\,\big|_{t=0} = -\int\rho\varphi(\log\rho + 1)$ along
  $\rho_t = \rho(1 + t\varphi)$ for bounded $\varphi$
  (\lean{MeasureTheory.hasDerivAt_integral_negMulLog_mul_one_add})
  \item \lean{ForMathlib/InformationTheory/MarginalEntropy.lean}:\newline the marginal
  $\rho_X(x) = \int\rho(x, y)\,dy$ of a density of finite entropy on $X \times Y$ has finite
  entropy if $\rho\log\gamma(y)$ and $\rho\log\gamma'(x)$ are integrable for positive $\gamma$,
  $\gamma'$ of integral at most~$1$
  (\lean{MeasureTheory.integrable_negMulLog_integral_prod}), and its moments are moments of
  $\rho$ (\lean{MeasureTheory.integral_mul_integral_prod})
  \item \lean{ForMathlib/InformationTheory/Subadditivity.lean}:\newline the entropy of a
  product $f(x)g(y)$ of probability densities is the sum of their entropies
  (\lean{MeasureTheory.integral_negMulLog_mul_prod}). If the marginals have finite entropy,
  the entropy is subadditive, $-\int\rho\log\rho \le -\int\rho_X\log\rho_X -
  \int\rho_Y\log\rho_Y$ (\lean{MeasureTheory.integral_negMulLog_le_add})
  \item \lean{ForMathlib/InformationTheory/ChangeOfVariables.lean}:\newline the density
  \lean{MeasureTheory.pushDensity} in new variables $u = \Phi(x)$, for $\Phi$ injective and
  differentiable with invertible derivative on a finite-dimensional space. Its moments are
  moments of $\rho$ (\lean{MeasureTheory.integral_mul_pushDensity}), and its entropy is
  $-\int\rho_\Phi\log\rho_\Phi = \int(-\rho\log\rho + \rho\log|\det\Phi'|)$
  (\lean{MeasureTheory.integral_negMulLog_pushDensity}).
  \end{itemize}
\end{itemize}

\paragraph{Setting.} An ensemble on a measure space $(X, \mu)$ (\lean{IsEnsemble}~$\rho$~$\mu$) is
a function $\rho \ge 0$ with $\int\rho\,d\mu = 1$ and $\rho\log\rho$ integrable (finite entropy).
Its entropy relative to a cell of size $c$ is \lean{entropy}~$c$~$\rho$~$\mu$ $= -\int
\rho\log(c\rho)\,d\mu$, in nats. A variable $f : X \to \mathbb R$ has mean \lean{mean}, variance
\lean{var} and standard deviation \lean{sd} under $\rho$. One DOF has phase space
$\mathbb R \times \mathbb R$ with Lebesgue measure and cell $h > 0$, and $\hbar = h/2\pi$
(\lean{hbar}). The conditions are:
\begin{itemize}
\item \lean{LowerEntropicBound}~$h$~$\rho$ (ZENT): $S_h[\rho] \ge 0$;
\item \lean{UncertaintyPrinciple}~$h$~$\rho$ (CUP): $\sigma_q\sigma_p \ge \hbar/e$.
\end{itemize}
Finite second moments of the variables are hypotheses of the theorems. The Poisson bracket
\lean{poissonBracket}~$f$~$g$ is $\partial_q f\,\partial_p g - \partial_p f\,\partial_q g$, with
Fréchet derivatives. For the brief's own proof,
\lean{lagrangian}~$h$~$\lambda_1$~$\lambda_2$~$S$~$\rho$ is
$\sigma_q^2\sigma_p^2 + \lambda_1(\int\rho - 1) + \lambda_2(S_h[\rho] - S)$, varied along
$\rho(1 + t\varphi)$ with $\varphi$ bounded and measurable, so $\delta\rho = \rho\varphi$, and
\lean{eulerLagrange} is the bracket of its first variation. New variables $(f, g)$
(\lean{IsNewVariables}~$f$~$g$) are differentiable, with $(q, p) \mapsto (f, g)$ injective and
$\{f, g\} \ne 0$; their density \lean{densityFG}~$f$~$g$~$\rho$ is $\rho/|\{f, g\}|$ at
$(f, g)(q, p)$ and $0$ off the range of $(f, g)$.

For several DOFs the phase space is $(\mathbb R^2)^{n+1}$ (\lean{Fin}~$(n + 1) \to \mathbb R
\times \mathbb R$), and DOF $i$ is the pair $(q_i, p_i)$ of coordinates $i$. Its marginal
\lean{marginal}~$\rho$~$i$ is $\rho$ integrated over the other $n$ DOFs. ZENT and CUP become
\lean{LowerEntropicBoundPhaseSpace}~$h$~$\rho$ ($S_h[\rho_i] \ge 0$ for every $i$) and
\lean{UncertaintyPrinciplePhaseSpace}~$h$~$\rho$ ($\sigma_{q_i}\sigma_{p_i} \ge \hbar/e$ for
every $i$, with $\sigma$ under $\rho$).

For the quantum case, $\rho_q, \rho_p$ are densities on $\mathbb R$, in quantum mechanics
$|\psi|^2$ and $|\hat\psi|^2$, with entropies $S_1$ and standard deviations $\sigma_q, \sigma_p$
of the identity. Their product \lean{pairDensity}~$\rho_q$~$\rho_p$ is a density on phase space.
The conditions are:
\begin{itemize}
\item \lean{EntropicUncertainty}~$h$~$\rho_q$~$\rho_p$, the bound of
  Bia\l{}ynicki-Birula and Mycielski: $S_1[\rho_q] + S_1[\rho_p] \ge \log(\pi e\hbar)$;
\item \lean{WehrlBound}~$h$~$Q$~$\rho_q$~$\rho_p$, Lieb's bound on the Wehrl entropy: for every
  $s > 0$, $Q_s$ is an ensemble on phase space with finite second moments,
  $\sigma^2_{Q_s,q} = \sigma_q^2 + s\hbar/2$, $\sigma^2_{Q_s,p} = \sigma_p^2 + \hbar/(2s)$ and
  $S_h[Q_s] \ge 1$. In quantum mechanics $Q_s$ is the Husimi function of $\psi$ for the coherent
  states squeezed by $s$, divided by $h$;
\item \lean{Heisenberg}~$h$~$\rho_q$~$\rho_p$: $\sigma_q\sigma_p \ge \hbar/2$.
\end{itemize}
For a classical ensemble $\rho$ on phase space, \lean{marginalFst}~$\rho$ and
\lean{marginalSnd}~$\rho$ are the densities $\int\rho(q, p)\,dp$ of $q$ and $\int\rho(q, p)\,dq$
of $p$.

\paragraph{Results.}
\begin{xltabular}{\textwidth}{@{}BL@{}}
\toprule
Brief & Lean \\
\midrule
\endfirsthead
\toprule
Brief & Lean \\
\midrule
\endhead
\bottomrule
\endlastfoot
Corollary~1, $S[\rho_U] = \log(\mu(U)/h^n)$ & \lean{entropy_uniform} (any measure space, any
  cell $c > 0$); \lean{entropy_uniform_phaseSpace} (phase space $(\mathbb R^2)^n$, $c = h^n$) \\
Corollary~1, $S[\rho_U] = 0$ iff $\mu(U) = h^n$ & \lean{entropy_uniform_eq_zero_iff} \\
ZENT, CUP & \lean{LowerEntropicBound}, \lean{UncertaintyPrinciple} (one DOF);
  \lean{LowerEntropicBoundPhaseSpace}, \lean{UncertaintyPrinciplePhaseSpace} (each DOF of
  $(\mathbb R^2)^{n+1}$) \\
Theorem~2, proof: a product of two independent Gaussians minimizes $\sigma_q\sigma_p$ at given
  entropy, and then $\sigma_q\sigma_p = (\hbar/e)e^{S}$ &
  \lean{hbar_div_exp_mul_exp_entropy_le}: $\sigma_q\sigma_p \ge (\hbar/e)e^{S_h[\rho]}$ for every
  ensemble; \lean{sd_mul_sd_eq_iff}: equality iff $\rho$ is a.e.\ a product of Gaussians \\
Theorem~2, eq.~(3), ZENT implies CUP & \lean{uncertaintyPrinciple_of_lowerEntropicBound};
  for several DOFs
  \lean{uncertaintyPrinciplePhaseSpace_of_lowerEntropicBoundPhaseSpace} \\
CUP, equality case & \lean{sd_mul_sd_eq_hbar_div_exp_iff}, corrected (see ``Problems found'');
  for several DOFs \lean{sd_mul_sd_eq_hbar_div_exp_iff_marginal} \\
ZENT, ``$\rho(q, p)$ is the marginal over the DOF'' & \lean{isEnsemble_marginal}: the marginal
  is an ensemble; \lean{mean_marginal}, \lean{var_marginal}: with the means and variances of
  $q_i, p_i$ under $\rho$ \\
Eq.~(4), the Poisson bracket is the Jacobian determinant & \lean{det_prod};
  \lean{poissonBracket_fst_snd}: $\{q, p\} = 1$ \\
Theorem~2, eq.~(5), $\sigma_f\sigma_g \ge (\hbar/e)e^{E[\log|\{f,g\}|]}$ under ZENT &
  \lean{hbar_div_exp_mul_exp_le_sd_mul_sd}; for every ensemble
  $\sigma_f\sigma_g \ge (\hbar/e)e^{S_h[\rho] + E[\log|\{f,g\}|]}$:
  \lean{hbar_div_exp_mul_exp_entropy_add_le}; by the brief's route through $\rho_{fg}$,
  \lean{IsNewVariables.hbar_div_exp_mul_exp_le_sd_mul_sd} \\
Eq.~(5), proof: $\rho_{qp} = \rho_{fg}|\{f, g\}|$ &
  \lean{IsNewVariables.densityFG_mul_abs_poissonBracket};
  \lean{IsNewVariables.isEnsemble_densityFG}: $\rho_{fg}$ is an ensemble \\
Eq.~(5), proof: $S[\rho_{fg}] = S[\rho_{qp}] + E[\ln|\{f, g\}|]$, and under ZENT
  $S[\rho_{fg}] \ge E[\ln|\{f, g\}|]$ & \lean{IsNewVariables.entropy_densityFG},
  \lean{IsNewVariables.integral_mul_log_le_entropy_densityFG} \\
Eq.~(5), proof: Gaussians minimize the uncertainty in $(f, g)$ &
  \lean{IsNewVariables.hbar_div_exp_mul_exp_entropy_densityFG_le}:
  $\sigma_f\sigma_g \ge (\hbar/e)e^{S_h[\rho_{fg}]}$; \lean{IsNewVariables.sd_mul_sd_eq_iff}:
  equality iff $\rho_{fg}$ is a.e.\ a product of Gaussians \\
Eq.~(5), equality case (not in the brief) &
  \lean{IsNewVariables.sd_mul_sd_eq_hbar_div_exp_mul_exp_iff}: under ZENT, equality iff
  $S_h[\rho] = 0$ and $\rho_{fg}$ is a.e.\ a product of Gaussians in $(f, g)$; then $(f, g)$ is
  onto up to a null set (\lean{volume_compl_range_eq_zero}) \\
Remark, CUP does not imply ZENT & \lean{exists_not_lowerEntropicBound} \\
Theorem~2, proof: $\delta L = \int\delta\rho\,[(q - \mu_q)^2\sigma_p^2 + \sigma_q^2(p - \mu_p)^2
  + \lambda_1 - \lambda_2\ln(h\rho) - \lambda_2]$ & \lean{hasDerivAt_lagrangian} \\
Theorem~2, proof: the solution of $\delta L = 0$ is a product of two independent Gaussians &
  \lean{hasDerivAt_lagrangian_zero_iff}: for $\rho > 0$ a.e., $\delta L = 0$ in every bounded
  direction iff $\rho$ is a.e.\ the product of Gaussians with the means and variances of
  $\rho$, $\lambda_2 = -2\sigma_q^2\sigma_p^2$ and
  $\lambda_1 = \lambda_2(1 + \ln(\hbar/\sigma_q\sigma_p))$ \\
Theorem~2, proof: then $\sigma_q\sigma_p = (\hbar/e)e^{S[\rho]}$ &
  \lean{sd_mul_sd_eq_of_hasDerivAt_lagrangian} \\
Conclusion: the quantum uncertainty principle follows from an entropic lower bound &
  Not with the von Neumann entropy (see ``Problems found''). From the bound of
  Bia\l{}ynicki-Birula and Mycielski: \lean{heisenberg_of_entropicUncertainty}, and
  \lean{sd_mul_sd_eq_hbar_div_two_iff}: equality iff $\rho_q, \rho_p$ are a.e.\ Gaussians with
  $\sigma_q\sigma_p = \hbar/2$. From Lieb's bound on the Wehrl entropy:
  \lean{heisenberg_of_wehrlBound}, and \lean{entropy_eq_one_of_wehrlBound}: at equality, at
  $s = \sigma_q/\sigma_p$, $S_h[Q_s] = 1$ and $Q_s$ is a.e.\ a product of Gaussians \\
(not in the brief) the classical counterpart &
  \lean{log_le_entropy_marginalFst_add_entropy_marginalSnd}: under ZENT,
  $S_1[\rho_q] + S_1[\rho_p] \ge \log h = \log(2\pi\hbar)$ for the marginals;
  \lean{log_pi_exp_hbar_sub_log}: the quantum bound exceeds it by $\log(e/2)$;
  \lean{uncertaintyPrinciple_of_lowerEntropicBound_marginal}: eq.~(3) by this route \\
(used for the quantum case) the Gaussian bound for two densities on $\mathbb R$ &
  \lean{exp_entropy_add_div_le_sd_mul_sd}:
  $\sigma_q\sigma_p \ge e^{S_1[\rho_q] + S_1[\rho_p]}/(2\pi e)$, the bound for one DOF applied to
  $\rho_q\rho_p$; \lean{sd_mul_sd_eq_exp_entropy_add_div_iff}: equality iff both are a.e.\
  Gaussians \\
(used for both equations) the Gaussian bound on any measure space &
  \lean{Hyp.hbar_div_exp_mul_exp_le}: $\sigma_f\sigma_g \ge (\hbar/e)e^{S_h[\rho] + E[\log J]}$
  for a weight $J > 0$ that turns every product of Gaussians in $(f, g)$ into a subprobability
  density (\lean{IsGaussWeight}) \\
\end{xltabular}

\paragraph{Differences from the brief.}
\begin{itemize}
\item The bound $\sigma_q\sigma_p \ge (\hbar/e)e^{S}$ is proved by Gibbs' inequality against
  the product $G$ of Gaussians in $q$ and $p$ with the means and variances of $\rho$:
  $-\int\rho\log\rho \le -\int\rho\log G = \log(2\pi e\,\sigma_q\sigma_p)$. This proves that the
  Gaussian is a minimum, not only a critical point, and gives the equality case. Only the
  marginal variances enter, so no covariance matrix is needed. The brief's Lagrange computation
  is formalized as well, as a second proof that the Gaussians are the critical points.
\item In the Lagrange computation, $\delta\rho$ is $\rho\varphi$ with $\varphi$ bounded and
  measurable, rather than arbitrary. Along these variations $\rho(1 + t\varphi)$ stays
  nonnegative for small $t$ and the entropy is differentiable. The vanishing of $\delta L$ in
  all of them gives the brief's equation for $\ln(h\rho)$ where $\rho > 0$. The brief solves it
  for $\rho$ in one line; here the cases $\lambda_2 = 0$ (the quadratic would be constant) and
  $\lambda_2 > 0$ ($\rho$ would be bounded below, so not integrable) are excluded, and
  normalization and the variances of $\rho$ fix the constant and $\lambda_2$.
\item For several DOFs, ZENT and CUP are stated on $(\mathbb R^2)^{n+1}$, so there is at least
  one DOF. An ``independent DOF'' is read as a pair $(q_i, p_i)$ of coordinates; DOFs in other
  canonical coordinates, after a change of variables that mixes the pairs, are not considered.
  The marginal of each DOF is shown to be an ensemble with finite second moments. Its finite
  entropy comes from Gibbs' inequality in the integrated variables, against a standard Gaussian
  in the other DOFs, and from the pointwise bound against a standard Gaussian in the DOF
  itself. Theorem~2 for one DOF then applies to each marginal. Corollary~1 is formalized for
  any $n$.
\item The hypotheses that the brief leaves implicit are explicit: $\rho \ge 0$, $\int\rho = 1$,
  finite entropy and finite second moments. Positive variances are proved, not assumed: an
  ensemble of finite entropy cannot have zero variance in any variable, since Gibbs' inequality
  against ever narrower Gaussians would bound its entropy by $-\infty$.
\item For eq.~(5) the brief passes to the density $\rho_{fg}$ of $(f, g)$, with
  $\rho_{qp} = \rho_{fg}\,|\{f, g\}|$. Both routes are formalized. In the first, $\rho_{fg}$ is
  not formed: the Gaussian in $(f, g)$ is pulled back to phase space with the weight
  $|\{f, g\}|$, and the change-of-variables formula shows that its integral is at most $1$. In
  the brief's, $\rho_{fg}$ is $\rho/|\{f, g\}|$ on the range of $(f, g)$ and $0$ off it, and the
  bound for one DOF is applied to it; this also gives the equality case. Both need $f, g$
  differentiable, $(q, p) \mapsto (f, g)$ injective (not surjective, not orientation
  preserving), $\{f, g\} \ne 0$ everywhere, and $\rho\log|\{f, g\}|$ integrable. $\sigma_f$ is
  the standard deviation of $f$ under $\rho$. Physical units are not modelled, so the brief's
  assumption on them is not needed.
\item The counterexample for the remark is a uniform distribution on two squares
  $[L, L + \varepsilon]^2$ and $[-L - \varepsilon, -L]^2$, with $\varepsilon = \sqrt h/2$ and
  $L = \sqrt h$. Its entropy is $\log\frac12 < 0$ and $\sigma_q\sigma_p \ge h/2 \ge \hbar/e$.
  The brief instead correlates $q$ and $p$.
\item For the quantum case, the brief points to the zero entropy of pure states. The
  formalization instead takes two entropic bounds of quantum mechanics as hypotheses: that of
  Bia\l{}ynicki-Birula and Mycielski (1975) on $S_1[|\psi|^2] + S_1[|\hat\psi|^2]$, and Lieb's
  (1978) on the Wehrl entropy, the entropy of the Husimi function. Both rest on Beckner's sharp
  Hausdorff--Young inequality, which is not in Mathlib, and are not proved here; nor are the
  Fourier transform and the Husimi function of a wave function formed. Each, with the brief's
  Gaussian bound, gives $\sigma_q\sigma_p \ge \hbar/2$:
  \begin{itemize}
  \item The bound of Bia\l{}ynicki-Birula and Mycielski is $S_1[\rho_q] + S_1[\rho_p] \ge
    \log(\pi e\hbar)$, and the Gaussian bound for one DOF, applied to $\rho_q\rho_p$, gives
    $\sigma_q\sigma_p \ge e^{S_1[\rho_q] + S_1[\rho_p]}/(2\pi e)$. Classically, ZENT and
    subadditivity give only $S_1[\rho_q] + S_1[\rho_p] \ge \log(2\pi\hbar)$, and the same bound
    gives CUP. The two principles differ by $\log(e/2)$ in this one inequality, the ratio of
    $\hbar/2$ to $\hbar/e$.
  \item Lieb's bound is $S_h[Q_s] \ge 1$, not $\ge 0$. With the variances of $Q_s$, the
    Gaussian bound gives $(\sigma_q^2 + s\hbar/2)(\sigma_p^2 + \hbar/(2s)) \ge \hbar^2$ for every
    $s > 0$, and at $s = \sigma_q/\sigma_p$ the left side is $(\sigma_q\sigma_p + \hbar/2)^2$.
    This argument is due to Grabowski (Rep.\ Math.\ Phys.\ 20, 153 (1984)); we follow the
    account in Ohya and Petz, \emph{Quantum Entropy and Its Use}, \S16, which goes through the
    marginals of $Q_s$ instead. With $S_h[Q_s] \ge 0$ the same steps give only
    $\sigma_q\sigma_p \ge (2/e - 1)\hbar/2 < 0$.
  \end{itemize}
  We checked the constants against Ohya--Petz and against the paper of Floerchinger, Haas and
  M\"uller-Groeling, arXiv:\allowbreak 2103.07229. The original papers of
  Bia\l{}ynicki-Birula--Mycielski, Lieb, Beckner and Grabowski could not be downloaded.
\end{itemize}

\paragraph{Problems found in the brief.} The main claim (ZENT implies CUP, eqs.~(3) and (5)) is
correct; the equality case is misstated, and the conclusion on the quantum case is not
supported.
\begin{itemize}
\item CUP says that equality holds ``if and only if position and momentum are independent
  gaussian distributions'', and the proof ends the same way. Independent Gaussians satisfy
  $\sigma_q\sigma_p = (\hbar/e)e^{S}$, which equals $\hbar/e$ only when $S = 0$. Under ZENT,
  $\sigma_q\sigma_p = \hbar/e$ holds exactly for products of Gaussians with
  $\sigma_q\sigma_p = \hbar/e$ (\lean{sd_mul_sd_eq_hbar_div_exp_iff}).
\item The Lagrange multipliers give only a stationary point of $\sigma_q^2\sigma_p^2$ at fixed
  entropy. They show neither that it is a minimum nor that a minimizer exists, which the proof
  then uses (``the gaussian is the distribution that minimizes the uncertainty''). The
  computation also needs $\rho > 0$ almost everywhere, which the brief does not state: a
  variation cannot move mass onto $\{\rho = 0\}$ while keeping the entropy differentiable,
  since $-x\ln x$ has infinite slope at $0$. So $\delta L = 0$ says nothing about where $\rho$
  vanishes (\lean{hasDerivAt_lagrangian_zero_iff} assumes $\rho > 0$ a.e.).
\item In the derivation of eq.~(5), ``order preserving'' should be ``orientation preserving'';
  the last two lines of the chain of equalities are identical; and CUP writes
  $E[\ln\{f, g\}]$ where the abstract and eq.~(5) have $E[\ln|\{f, g\}|]$. The logarithm needs
  the absolute value when $\{f, g\} < 0$.
\item The brief does not state that the entropy and the second moments must be finite. Without
  this $S[\rho]$ or $\sigma$ may be undefined. For several DOFs, finite entropy of $\rho$ does
  not even give its marginals finite entropy (not formalized): for any probability density
  $\rho_1$ on $\mathbb R$, the indicator of $\{(x, y) : |y| < \rho_1(x)/2\}$ has entropy $0$
  (for $h = 1$) and marginal $\rho_1$, whose entropy can be $-\infty$. Here the second moment
  of $y$ is infinite; finite second moments of all DOFs exclude this
  (\lean{isEnsemble_marginal}).
\item The conclusion says that the lower bound on the entropy ``is already present in quantum
  mechanics as pure states have zero entropy'', and that ``in general, the uncertainty principle
  is a consequence of the entropic lower bound''. The von Neumann entropy is $\ge 0$ for every
  state and $= 0$ for every pure state, whatever its $\sigma_q\sigma_p$, so this bound
  constrains nothing; Heisenberg's principle follows from $[q, p] = i\hbar$ alone (Robertson; in
  physlib \lean{LinearPMap.state_uncertainty_of_raw_commutator}).
  A positive bound is needed even classically, since the brief's argument turns $S \ge 0$ into
  $\hbar/e < \hbar/2$. The principle does follow from the two quantum entropic bounds above,
  which are stronger than zero entropy: by $\log(e/2)$ for the marginals, and with $1$ in place of
  $0$ for the Wehrl entropy.
\end{itemize}

\end{document}
